\documentstyle[% 


righttag% 


,11pt% 


,amssymb% 


,russian%               remove in Eng. 


,izvru%                 Set izven% in Eng. 


% ,izvru-ks             for brief communications 


]{amsart} 


 


%       Math definitions 


 


%\newcommand{\wt}{\widetilde} 


%\newcommand{\ov}{\overline} 


 


%\hoffset=-15mm                  For Eng. 


 


\setcounter{page}{13} 


\god{1997} 


\nomer{9~(424)} %                Remove in Eng. 


%\nomer{Vol.\,41, No.\,9}        Set in Eng. 


% \str{75-81}                    Set in Eng. 


\udk{514.755} 


\author{\it Т.Б.\,Жогова} 


\title{Проективное изгибание второго порядка семейств $L^m_{2n-1}$} 


 


\date{} 


% \pagestyle{myheadings}         Set in Eng. 


 


\pagestyle{plain} %              Remove in Eng. 


\begin{document} 


 


% \thispagestyle{plain}          Set in Eng. 


 


\maketitle 


 


% Body text 


 


Условимся, что по индексам $i$, $j$ суммирования нет, а по индексам $p$, 


$q$, 


$r$ всегда проводится суммирование. Все различные индексы, по которым не 


проводится суммирование, находящиеся в одном и том же математическом 


выражении, не принимают равные значения. Все индексы, как правило, 


принимают значения от 1 до $n$ включительно. 


 


Будем рассматривать классическое изгибание Фубини-Картана (\cite{Fin}, 


с.\,102) погруженного многообразия $L^m_{2n-1}$ в проективное пространство 


$P_{2n-1}$. 


 


Изгибание первого порядка семейств $L^m_{2n-1}$ исследовано автором в 


работе \cite{Z1}. 


При $n=3$, $m=2$ задача проективного изгибания 2-го порядка рассмотрена 


автором в работе \cite{Z2}. 


 


В $P_{2n-1}$ введем проективный репер, состоящий из $2n$ линейно 


независимых 


аналитических точек $A_i$, $A_{n+j}$, инфинитезимальные перемещения 


которого определим системой дифференциальных уравнений 


$$ 


dA_i=\omega^p_iA_p+\omega^{n+p}_iA_{n+p},\quad 


dA_{n+i}=\omega^p_{n+i}A_p+\omega^{n+p}_{n+i}A_{n+p}, 


$$ 


в которой формы Пфаффа удовлетворяют известным уравнениям структуры 


проективного пространства. 


 


Пусть $(n-1)$-плоскость 


$$ 


L_{n-1}=(A_1,\dotsc,A_n) 


$$ 


описывает семейство $L^m_{2n-1}$. Отнесем семейство $L^m_{2n-1}$ к реперу 


1-го порядка: точки $A_i$ 


поместим в фокусы $(n-1)$-плоскости $L_{n-1}$, фокальное направление 


фокуса $A_i$ определим 


уравнением $\omega_i^{n+i}=0$ и точку $A_{n+i}$ поместим на прямую Лапласа 


фокуса $A_i$. Тогда семейство $L^m_{2n-1}$ определится системой уравнений 


Пфаффа 


\begin{equation} 


\omega_i^{n+j}=0. 


\end{equation} 


Внешним дифференцированием уравнений (1) находим 


\begin{equation} 


\omega^j_i\wedge\omega_j-\omega^{n+j}_{n+i}\wedge\omega_i=0, 


\end{equation} 


где $\omega_i=\omega_i^{n+i}$. За независимые формы семейства $L^m_{2n-1}$ 


принимаются формы $\omega_\alpha$ 


($\alpha=\ov{1,m}$, $2\le m\le n$). 


 


Рассмотрим семейство $\ov L^m_{2n-1}$, которое изгибанием $k$-го порядка 


наложимо на семейство $L^m_{2n-1}$. Отнесем семейство 


$\ov L^m_{2n-1}$ к проективному 


реперу $\{B_i,B_{n+j}\}$, 


инфинитезимальные перемещения которого определим системой дифференциальных 


уравнений 


$$ 


dB_i=\Omega^p_iB_p+\Omega_i^{n+p}B_{n+p},\quad 


dB_{n+i}=\Omega^p_{n+i}B_p+\Omega^{n+p}_{n+i}B_{n+p}. 


$$ 


 


Из работы \cite{Z1} следует, что без ограничения общности семейство 


$\ov L^m_{2n-1}$ так 


же, как 


и семейство $L^m_{2n-1}$, можно отнести к реперу 1-го порядка. Тогда будут 


иметь место уравнения 


\begin{equation} 


\Omega^{n+j}_i=0. 


\end{equation} 


 


Рассмотрим проективное преобразование $\Pi$, переводящее репер 


$\{B_i,B_{n+j}\}$ в $\{\wt B_i,\wt B_{n+j}\}$, совпадающий с 


$\{A_i,A_{n+j}\}$, так что $\wt B_i=A_i$, $\wt B_{n+j}=A_{n+j}$. 


 


Условие касания семейств 


$$ 


\wt L^m_{2n-1}=\Pi(\ov L^m_{2n-1})\quad\text{и}\quad L^m_{2n-1} 


$$ 


2-го порядка определяется уравнением 


\begin{multline} 


(\wt B_1,\dotsc,\wt B_n)+d(\wt B_1,\dotsc,\wt B_n)+ 


\frac{1}{2}d^2(\wt B_1,\dotsc,\wt B_n)= \\


=(\theta_0+\theta_1+\theta_2)\{(A_1,\dotsc,A_n) 


+d(A_1,\dotsc,A_n)+\frac{1}{2}d^2(A_1,\dotsc,A_n)\}, 


\end{multline} 


где $\theta_k$ --- некоторые скалярные множители, и индекс $k$ указывает на 


порядок малости такого множителя относительно независимых форм 


$\omega_\alpha$. 


 


Введем обозначения 


$$ 


\Omega^{n+i}_i=\Omega_i, \quad 


\Tilde\omega^\delta_\gamma=\Omega^\delta_\gamma-\omega^\delta_\gamma 


\quad(\delta,\gamma=\ov{1,2n}). 


$$ 


Условие (4) в силу уравнений (1), (3) приводит к следующей системе 


уравнений:


\begin{gather} 


\theta_0=1,\quad \theta_1=\Tilde\omega^r_r,\quad 


\theta_2=d\Tilde\omega^r_r+\Tilde\omega^r_r\Tilde\omega^p_p+ 


\omega_q\Tilde\omega^q_{n+q}\notag,\\ 


\Tilde\omega^j_i\omega_j-\Tilde\omega^{n+j}_{n+i}\omega_i=0,\\ 


\Tilde\omega_i=0,\quad \Tilde\omega^i_i-\Tilde\omega^{n+i}_{n+i}=0. 


\end{gather} 


Дифференцируя внешним образом уравнения (3), получим 


\begin{equation} 


\Tilde\omega^j_i\wedge\omega_j-\Tilde\omega^{n+j}_{n+i} 


\wedge\omega_i=0. 


\end{equation} 


Из уравнений (5) и (7) следует


\begin{equation} 


\Tilde\omega^j_i=0,\quad \Tilde\omega^{n+j}_{n+i}=0. 


\end{equation} 


Внешним дифференцированием уравнений (6), (8) находим 


\begin{equation} 


\Tilde\omega^i_{n+i}\wedge\omega_i=0, 


\end{equation} 


\begin{equation} 


\begin{split} 


(\Tilde\omega^i_i-\Tilde\omega^j_j)\wedge\omega^j_i- 


\Tilde\omega^j_{n+i}\wedge\omega_i&=0,\\ 


(\Tilde\omega^i_i-\Tilde\omega^j_j)\wedge\omega^{n+j}_{n+i}+ 


\Tilde\omega^j_{n+i}\wedge\omega_j&=0.


\end{split} 


\end{equation} 


Замкнутая система уравнений (1), (2), (6), (8), (9), (10) определяет 


семейство $L^m_{2n-1}$ вместе со своим изгибанием 2-го порядка. 


Исследование этой 


системы показывает, что при $m=2$ она находится в инволюции с характерами 


$s_1=3n^2-2n$, $s_2=2n-3$, т.\,е. справедлива 


 


\begin{thm} Семейство $L^2_{2n-1}$, допускающее проективное изгибание 


{\em 2}-го 


порядка, существует с произволом $2n-3$ функций двух аргументов. 


\end{thm} 


 


Если же $m=n$, то рассматриваемая замкнутая система не находится в 


инволюции, 


и ее необходимо продолжать. Сделаем частичное продолжение этой системы. 


 


Раскрывая по лемме Картана уравнения (2), получим 


\begin{equation} 


\omega^j_i=a^j_i\omega_j+c^j_i\omega_i,\quad 


\omega^{n+j}_{n+i}=b^j_i\omega_i-c^j_i\omega_j. 


\end{equation} 


Из системы (10) в силу (11) следует, что форма 


$\Tilde\omega^i_i-\Tilde\omega^1_1$ разлагается только по 


формам $\omega_i$ и $\omega_1$: 


$$ 


\Tilde \omega^i_i-\Tilde\omega^1_1=a_i\omega_i+b_i\omega_1. 


$$ 


Отсюда находим


$$ 


\Tilde\omega^j_j-\Tilde\omega^1_1=a_j\omega_j+b_j\omega_1 


$$ 


и, следовательно, 


$$ 


\Tilde\omega^i_i-\Tilde\omega^j_j= 


a_i\omega_i-a_j\omega_j+(b_i-b_j)\omega_1. 


$$ 


Отсюда вытекает, что $b_i=b_j=a_1$, т.\,е. 


\begin{equation} 


\Tilde\omega^i_i-\Tilde\omega^1_1=a_i\omega_i+a_1\omega_1. 


\end{equation} 


Продифференцировав внешним образом эти уравнения, получим 


$$ 


\Delta a_i\wedge\omega_i+\Delta a_1\wedge\omega_1=0, 


$$ 


где 


$$ 


\Delta a_j=da_j+a_j(\omega^j_j-\omega^{n+j}_{n+j}). 


$$ 


Наиболее общее алгебраическое решение этой системы квадратичных уравнений 


имеет вид 


$$ 


\Delta a_j=p_j\omega_j 


$$ 


и, следовательно, 


\begin{equation} 


\Delta a_j\wedge\omega_j=0. 


\end{equation} 


 


Заметим, что сделанное частичное продолжение рассматриваемой системы (при 


помощи уравнений (12)) в силу теоремы 2.2 \cite{Mak1} является 


корректным. 


 


Система уравнений (10) с учетом (12) принимает вид 


\begin{equation} 


\begin{split} 


\omega^j_i\wedge(a_i\omega_i-a_j\omega_j) 


+\Tilde\omega^j_{n+i}\wedge\omega_i&=0,\\ 


\omega^{n+j}_{n+i}\wedge(a_i\omega_i-a_j\omega_j) 


-\Tilde\omega^j_{n+i}\wedge\omega_j&=0.


\end{split} 


\end{equation} 


 


Замкнутая система уравнений (1), (3), (6), (8), (12); (2), (9), (13), (14) 


определяет семейство $L^n_{2n-1}$ вместе со своим изгибанием 2-го порядка. 


Эта система содержит $q=n(3n-1)$ искомых форм


$$ 


\omega^j_i,\quad \omega^{n+j}_{n+i},\quad 


\Tilde\omega^j_{n+i},\quad \Tilde \omega^i_{n+i}, 


\quad\Delta a_j. 


$$ 


Матрица полного продолжения этой системы является $m$-блочной \cite{Mak2}. 


Не равные 


нулю характеры уравнений (9) и (13) при фиксированных $i$ и $j$ 


соответственно равны 


$$ 


s^i_1=1,\quad s^j_1=1\quad\text{и}\quad N_i=N_j=1. 


$$ 


Таких блоков $2n$. Остальные блоки матрицы полного продолжения 


соответствуют 


системе ковариантов (2), (14) для каждой пары фиксированных индексов 


$i$ и $j$. Таких блоков $n(n-1)$. 


 


Положим 


$$ 


\Tilde\omega^j_{n+i}=p^j_i\omega_i+q^j_i\omega_j. 


$$ 


В силу (11) система уравнений (2), (14) приводит к конечным соотношениям 


$$ 


p^j_i=a_ic^j_i-a_jb^j_i,\quad 


q^j_i=-a_ia^j_i-a_jc^j_i, 


$$ 


т.\,е. $N_{ij}=3$: $a^j_i$, $b^j_i$, $c^j_i$. Ранг полярной системы 


уравнений (2), (14) равен $s^{ij}_1=3$, если 


$$ 


\omega_i\omega_j(a_i\omega_i-a_j\omega_j)\ne 0,\quad 


\text{т.\,е.}\quad N_{ij}=Q_{ij}. 


$$ 


 


Таким образом, исследуемая система уравнений находится в инволюции с 


характером 


$$ 


s_1=ns^i_1+ns^j_1+n(n-1)s^{ij}_1= 


2n+3n(n-1)=3n^2-n, 


$$ 


т.\,е. справедлива 


 


\begin{thm} Семейство $L^n_{2n-1}$, допускающее проективное изгибание 


{\em 2}-го 


порядка, существует с произволом $3n^2-n$ функций одного аргумента. 


\end{thm} 


 


Теперь следует остановиться на статье Л.\,М.Шепеленко \cite{Shep}, в 


которой автор рассматривает проективное изгибание семейств плоскостей 


Коровина-Гейдельмана, т.\,е. семейств $L^n_{2n-1}$. При исследовании 


системы, 


определяющей изгибание 2-го порядка, автором допущена ошибка при 


определении произвола наиболее общего интегрального элемента. В статье 


утверждается, что 


$$ 


N=3p^2-2,\quad s_1=3p^2-2p,\quad s_2=p-1, 


$$ 


в то время как $N=3p^2-2p$, т.\,е. система не в инволюции и, 


следовательно, 


теорема, сформулированная в статье (\cite{Shep}, с.\,35), не справедлива. 


 


Наконец, рассмотрим семейство $L^m_{2n-1}$, для которого 


$m\ne 2,n$. В этом случае 


формы $\omega_\beta$ 


($\beta=\ov{m+1,n}$) будут являться искомыми и, следовательно, 


\begin{equation} 


\omega_\beta=\Lambda^\alpha_\beta\omega_\alpha\quad 


(\alpha=\ov{1,m}). 


\end{equation} 


Продифференцировав внешним образом эти уравнения, получим 


\begin{equation} 


\Delta\Lambda^\alpha_\beta\wedge\omega_\alpha=0, 


\end{equation} 


где 


$$ 


\Delta\Lambda^\alpha_\beta=d\Lambda^\alpha_\beta+ 


\Lambda^\alpha_\beta(\omega^\alpha_\alpha-\omega^\beta_\beta- 


\omega^{n+\alpha}_{n+\alpha}+\omega^{n+\beta}_{n+\beta})\quad 


\text{(не суммировать)}. 


$$ 


 


Замкнутая система уравнений (1), (2), (15), (16) определяет семейство 


$L^m_{2n-1}$. Присоединяя к замкнутой системе уравнений (1), (3), (6), 


(8), (12); 


(2), (9), (13), (14), (система 1), определяющей семейство $L^n_{2n-1}$ 


вместе 


со своим изгибанием 2-го порядка, систему (15), (16), получим замкнутую 


систему уравнений, которая определяет семейство $L^m_{2n-1}$ вместе со 


своим изгибанием 2-го порядка. 


 


Присоединяемая система уравнений (15), (16) не нарушает инволютивности 


системы 1, а замкнутая система уравнений (15), (16) находится в инволюции 


с характерами $s_1=\dotsb=s_m=n-m$, т.\,е. справедлива 


 


\begin{thm} При $m\ne 2,n$ существуют с произволом $n-m$ функций $m$ 


аргументов семейства $L^m_{2n-1}$, допускающие проективное изгибание 


{\em 2}-го порядка. 


\end{thm} 


 


В заключение отметим, что проективное изгибание 3-го порядка семейств 


$L^m_{2n-1}$ будет тривиальным, т.\,е. семейства $L^m_{2n-1}$ и 


$\ov L^m_{2n-1}$ являются 


проективно эквивалентными. 


 


\begin{thebibliography}{9} 


 


\bibitem{Fin} 


Фиников С.П. {\em Теория пар конгруэнций}. -- М.: Гостехиздат, 1956. -- 


443 с. 


\bibitem{Z1} 


Жогова Т.Б. {\em О проективном изгибании семейств $L^m_{2n-1}$} // Учредит.


конф. Российск. ассоц. ``Женщины-математики'': Тез. докл. -- Москва--Суздаль, 


1993. -- С.\,20. 


\bibitem{Z2} 


Жогова Т.Б. {\em К вопросу о проективном изгибании двупараметрических 


семейств двумерных плоскостей в $P^5$}. -- Ред. журн. `` Изв. вузов. 


Математика''. -- Казань, 1979. -- 16~с. -- Деп. в ВИНИТИ 23.07.79, 


\No\,2761--79. 


\bibitem{Mak1} 


Макеев Г.Н. {\em К вопросу об инволютивности систем уравнений Пфаффа} // 


Изв. вузов. Математика. -- 1980. -- \No\,1. -- С.\,39--44. 


\bibitem{Mak2} 


Макеев Г.Н. {\em О некоторых признаках инволютивности систем уравнений 


Пфаффа} // Изв. вузов. Математика. -- 1982. -- \No\,9. -- С.\,81--83. 


\bibitem{Shep} 


Шепеленко Л.М. {\em Проективное изгибание расслояемых семейств плоскостей 


Коровина-Гейдельмана} // Геометрич. сб. -- Томск, 1967. -- Вып.\,6. -- 


С.\,32--35. 


 


 


\end{thebibliography} 


 


\vskip 1cm 


 


\post{\it Нижегородский государственный}{\it Поступила} 


\post{\it педагогический университет}{06.02.1995} 


 


\end{document} 


